\documentclass[11pt,a4paper]{article}

\usepackage[T1]{fontenc}
\usepackage[utf8]{inputenc}
\usepackage{lmodern}
\usepackage{microtype}
\usepackage[margin=2.5cm,headheight=14pt]{geometry}
\usepackage[english]{babel}
\usepackage[style=authoryear,maxcitenames=2,backend=biber]{biblatex}
\usepackage{imakeidx}
\usepackage[hidelinks]{hyperref}
\usepackage{twaidict}

\addbibresource{references.bib}
\makeindex[intoc,columns=2,title=Index of Terms]
\indexsetup{othercode=\small}

\hypersetup{
  pdftitle   = {TwAI Dictionary of AI, ML, and Trustworthy AI Terms},
  pdfsubject = {Terminology with multiple sourced definitions},
}

\title{TwAI Dictionary\\[4pt]
  \large Artificial Intelligence, Machine Learning, and Trustworthy AI Terms}
\author{TwAI Guidelines}
\date{\today}

\begin{document}

\pagestyle{empty}
\maketitle
\thispagestyle{empty}

\section*{How to Read This Dictionary}

Terms in artificial intelligence are defined differently by standards bodies,
regulators, and researchers. The differences matter: a system can be an
\enquote{AI system} under one definition and not another, and
\enquote{fairness} names several mathematically incompatible criteria. This
dictionary therefore gives \emph{several} definitions for each term, each tied
to its source, rather than a single preferred one.

\begin{itemize}[itemsep=2pt]
  \item Definitions in \enquote{quotation marks} are quoted from the source.
  \item Definitions marked \emph{adapted from} are paraphrased or condensed.
        Consult the source before citing them as formal text.
  \item A \emph{Note} after the definitions explains how they differ, and
        \emph{See also} links to related entries.
  \item Terms are grouped into three domains: general artificial intelligence,
        machine learning, and trustworthy AI. Within each domain they are
        ordered alphabetically. They are followed by Useful Notes, which
        give guidance spanning several terms, and a list of acronyms. An
        index of all terms and abbreviations is at the end.
\end{itemize}

\tableofcontents

\clearpage
\pagestyle{fancy}

% General artificial intelligence terms, in alphabetical order.
\domain{Artificial Intelligence}

\begin{entry}{agent}{Agent}
  \defn*{russell-norvig-2021}{An agent is anything that can be viewed
    as perceiving its environment through sensors and acting upon that
    environment through actuators.}
  \defn{wooldridge-1995}{A hardware or, more usually, software-based computer
    system with four properties: autonomy (it operates without direct human
    intervention), social ability (it interacts with other agents or humans),
    reactivity (it perceives and responds to changes in its environment), and
    pro-activeness (it takes the initiative to pursue goals).}
  \defn{iso-22989}{An automated entity that senses and responds to its
    environment and takes actions to achieve its goals.}
  \note{Russell and Norvig's definition is deliberately broad and includes a
    thermostat. Wooldridge and Jennings add requirements, such as autonomy and
    goal-directedness, that separate agents from ordinary programs.}
  \related{agentic-ai,autonomy,ai-system}
\end{entry}

\begin{entry}{agentic-ai}{Agentic AI}[AI agent]
  \defn{shavit-2023}{AI systems that can pursue complex goals with limited
    direct supervision, including by taking actions in the world that were
    not specified step by step by a human.}
  \defn{chan-2023}{Algorithmic systems that show increasing degrees of
    agency. Agency has four characteristics: goals are given without
    specifying how to reach them; actions have an impact without a human
    intermediary; behaviour is goal-directed; and the system plans over the
    long term.}
  \defn{owasp-llm-2025}{LLM-based systems granted \emph{agency}: the ability
    to call functions, use tools, or interface with other systems in
    response to prompts. \enquote{Excessive agency} is the risk that such
    functionality, permissions, or autonomy go beyond what the task needs.}
  \note{\enquote{Agentic} is best understood as a matter of degree rather
    than a binary property. Governance frameworks usually focus on the
    delegation of actions, and on who is accountable for them.}
  \related{agent,autonomy,human-oversight,prompt-injection}
\end{entry}

\begin{entry}{agi}{Artificial General Intelligence}[AGI]
  \defn{goertzel-2007}{AI systems with the ability to solve a wide variety of
    problems in a wide variety of contexts, including problems they were not
    specifically designed for, as opposed to systems built for one narrow
    task.}
  \defn{morris-2024}{A system whose capabilities match or exceed those of
    skilled humans across a wide range of non-physical tasks, including the
    ability to learn new tasks. AGI is best characterised by levels of
    performance and generality, not a single threshold.}
  \defn{searle-1980}{Related to \enquote{strong AI}: the claim that an
    appropriately programmed computer does not merely simulate a mind but
    literally has cognitive states and understanding.}
  \note{No agreed operational definition exists. Searle's strong AI is a
    philosophical claim about minds, whereas modern usage, as in Morris et al.,
    is about measurable capability.}
  \related{narrow-ai,artificial-intelligence}
\end{entry}

\begin{entry}{artificial-intelligence}{Artificial Intelligence}[AI]
  \defn*{mccarthy-2007}{It is the science and engineering of making
    intelligent machines, especially intelligent computer programs.}
  \defn*{nilsson-2010}{Artificial intelligence is that activity
    devoted to making machines intelligent, and intelligence is that quality
    that enables an entity to function appropriately and with foresight in
    its environment.}
  \defn{russell-norvig-2021}{The study and construction of agents that act
    rationally, meaning that they do the right thing given what they know,
    as opposed to the study of systems that think or act like humans.}
  \defn{iso-22989}{The research and development of mechanisms and
    applications of AI systems. The term also names the discipline itself.}
  \note{The earlier academic definitions describe a \emph{field of study},
    while standards and regulation mostly define the \emph{AI system}. When
    applicability or compliance is at stake, use the system-level definitions.}
  \related{ai-system,machine-learning,agent}
\end{entry}

\begin{entry}{ai-system}{AI System}
  \defn*[Art.~3(1)]{eu-ai-act}{A machine-based system that is designed to
    operate with varying levels of autonomy and that may exhibit
    adaptiveness after deployment, and that, for explicit or implicit
    objectives, infers, from the input it receives, how to generate outputs
    such as predictions, content, recommendations, or decisions that can
    influence physical or virtual environments.}
  \defn*{oecd-2024}{A machine-based system that, for explicit or implicit
    objectives, infers, from the input it receives, how to generate outputs
    such as predictions, content, recommendations, or decisions that can
    influence physical or virtual environments. Different AI systems vary in
    their levels of autonomy and adaptiveness after deployment.}
  \defn*{nist-ai-rmf}{An engineered or machine-based system that can,
    for a given set of objectives, generate outputs such as predictions,
    recommendations, or decisions influencing real or virtual environments.}
  \defn{iso-22989}{An engineered system that generates outputs such as
    content, forecasts, recommendations, or decisions for a given set of
    human-defined objectives.}
  \note{The EU and OECD definitions were aligned in 2023–2024 and both make
    \emph{inference} the distinguishing feature. That excludes conventional
    software whose rules are defined only by humans. The NIST and ISO/IEC
    definitions predate this change and do not use the inference test.}
  \related{artificial-intelligence,autonomy,general-purpose-ai}
\end{entry}

\begin{entry}{autonomy}{Autonomy}
  \defn[Recital~12]{eu-ai-act}{The degree to which an AI system's actions
    are independent of human involvement, and its ability to operate without
    human intervention.}
  \defn{iso-22989}{A characteristic of a system that can modify its intended
    domain of use or goals without external intervention, control, or
    oversight. Levels of automation run from none, through heteronomous
    systems under human control, to autonomous systems.}
  \defn{wooldridge-1995}{The ability of an agent to operate without the
    direct intervention of humans or others, having some control over its
    own actions and internal state.}
  \note{ISO/IEC 22989 separates \emph{automation}, which is carrying out a
    human-defined task, from \emph{autonomy}, which includes changing goals.
    Much everyday use of \enquote{autonomous} refers only to high automation.}
  \related{agent,agentic-ai,human-oversight}
\end{entry}

\begin{entry}{foundation-model}{Foundation Model}
  \defn*{bommasani-2021}{Any model that is trained on broad data
    (generally using self-supervision at scale) that can be adapted (e.g.,
    fine-tuned) to a wide range of downstream tasks.}
  \defn{nist-ai-600-1}{Discussed via the \enquote{dual-use foundation
    model}: an AI model trained on broad data, generally using
    self-supervision, with at least tens of billions of parameters, that is
    applicable across a wide range of contexts and could pose serious risks
    to security, public health, or safety.}
  \defn[Art.~3(63)]{eu-ai-act}{Treated in the regulation as a
    \enquote{general-purpose AI model}. See that entry.}
  \note{The term was coined to stress that one model is a shared, and
    therefore single-point-of-failure, basis for many applications.
    Regulators generally use \enquote{general-purpose AI} instead.}
  \related{general-purpose-ai,large-language-model,fine-tuning,self-supervised-learning}
\end{entry}

\begin{entry}{general-purpose-ai}{General-Purpose AI Model}[GPAI]
  \defn[Art.~3(63)]{eu-ai-act}{An AI model, including one trained with a
    large amount of data using self-supervision at scale, that displays
    significant generality, is capable of competently performing a wide range
    of distinct tasks, and can be integrated into a variety of downstream
    systems or applications. Models used only for research, development, or
    prototyping before being placed on the market are excluded.}
  \defn[Art.~3(66)]{eu-ai-act}{A \emph{general-purpose AI system} is an AI
    system based on a general-purpose AI model that can serve a variety of
    purposes, both for direct use and for integration into other AI systems.}
  \note{The EU AI Act separates the \emph{model} from the \emph{system} built
    on it and assigns obligations to each. It adds further obligations for
    models with \enquote{systemic risk}.}
  \related{foundation-model,ai-system}
\end{entry}

\begin{entry}{generative-ai}{Generative AI}[GenAI]
  \defn{nist-ai-600-1}{The class of AI models that emulate the structure and
    characteristics of input data in order to generate derived synthetic
    content, such as images, videos, audio, text, and other digital content.}
  \defn[ch.~20]{goodfellow-2016}{Models that learn a probability
    distribution over data, explicitly or implicitly, and can draw new
    samples from it. Examples are variational autoencoders, generative
    adversarial networks, and autoregressive models.}
  \defn[Art.~50]{eu-ai-act}{AI systems, including general-purpose AI
    systems, that generate synthetic audio, image, video, or text content.
    Their providers must mark outputs as artificially generated in a
    machine-readable format.}
  \note{The policy definitions refer to outputs (synthetic content). The
    technical definition refers to what the model learns (a data
    distribution).}
  \related{large-language-model,foundation-model,hallucination}
\end{entry}

\begin{entry}{hallucination}{Hallucination}[confabulation]
  \defn{ji-2023-hallucination}{Generated content that is nonsensical or
    unfaithful to the provided source content. \emph{Intrinsic}
    hallucinations contradict the source, and \emph{extrinsic}
    hallucinations cannot be verified from it.}
  \defn{nist-ai-600-1}{Termed \emph{confabulation}: the production of
    confidently stated but erroneous or false content, which may mislead or
    deceive users.}
  \defn{zhao-2023}{Content generated by a large language model that is
    factually incorrect or inconsistent with its input or context, but
    presented fluently and plausibly.}
  \note{NIST prefers \enquote{confabulation} to avoid attributing perception
    to the system. Research usage separates unfaithfulness to a given source
    from factual error about the world.}
  \related{large-language-model,retrieval-augmented-generation,validity-reliability}
\end{entry}

\begin{entry}{large-language-model}{Large Language Model}[LLM]
  \defn{zhao-2023}{A Transformer-based language model with a very large
    number of parameters, typically billions or more, that is pre-trained on
    massive text corpora and shows abilities, such as in-context learning,
    that are not found in smaller models.}
  \defn{brown-2020}{An autoregressive language model that, at sufficient
    scale, can perform new tasks from a natural-language instruction or a few
    examples in its prompt, without updating its weights.}
  \defn{bommasani-2021}{The most prominent kind of foundation model: a
    language model trained on broad text data that is adapted, by
    fine-tuning or prompting, to many downstream language tasks.}
  \note{\enquote{Large} is relative and has shifted over time. The defining
    features in practice are pre-training at scale and general-purpose
    prompting.}
  \related{foundation-model,transformer,prompt,hallucination}
\end{entry}

\begin{entry}{narrow-ai}{Narrow AI}[weak AI]
  \defn{goertzel-2007}{AI systems designed and trained for a specific task or
    a narrow range of tasks, which cannot generalise outside that domain
    without being redesigned or retrained.}
  \defn{searle-1980}{\enquote{Weak AI}: the position that computers are
    powerful tools for studying and simulating the mind, without any claim
    that the simulation itself has understanding.}
  \defn{morris-2024}{Systems with high performance on a clearly scoped task,
    such as a chess engine or a protein-structure predictor, placed on the
    \emph{narrow} end of a generality scale.}
  \related{agi,artificial-intelligence}
\end{entry}

\begin{entry}{prompt}{Prompt}
  \defn{liu-2023-prompt}{The input text, often a template with slots to be
    filled, that is given to a pre-trained language model so that the
    desired task is framed as text the model can continue or complete.}
  \defn{nist-ai-100-2}{The input to a generative AI model. In deployed
    applications it combines a \emph{system prompt} written by the
    application developer with user input and, often, content retrieved from
    external resources.}
  \defn{brown-2020}{The context given to a model at inference time, which may
    contain a task description and zero, one, or a few demonstrations, as in
    \emph{zero-shot}, \emph{one-shot}, and \emph{few-shot} prompting.}
  \related{large-language-model,prompt-injection}
\end{entry}

\begin{entry}{retrieval-augmented-generation}{Retrieval-Augmented Generation}[RAG]
  \defn{lewis-2020}{A method that combines a pre-trained parametric
    generator with a non-parametric memory, such as a dense vector index of
    documents. Retrieved passages condition the generation.}
  \defn{gao-2023-rag}{A paradigm that enhances large language models by
    retrieving relevant information from external knowledge bases at
    inference time. It improves accuracy and credibility, especially for
    knowledge-intensive tasks, and allows knowledge to be updated without
    retraining.}
  \defn{zhao-2023}{A form of \emph{retrieval augmentation} for large language
    models that injects relevant external knowledge into the prompt to
    improve factuality and to give access to up-to-date or private
    information.}
  \note{RAG reduces hallucination but introduces new risks. Retrieved
    content can carry an indirect prompt injection.}
  \related{hallucination,prompt-injection,embedding}
\end{entry}

% Machine learning terms, in alphabetical order.
\domain{Machine Learning}

\begin{entry}{datasets}{Training, Validation, and Test Data}
  \defn{iso-22989}{\emph{Training data} are samples used to train a machine
    learning model. \emph{Validation data} are samples used to tune
    hyperparameters or compare candidate models. \emph{Test data} are
    samples, independent of the training data, used to assess the performance
    of the final model.}
  \defn[§7.2]{hastie-2009}{In a data-rich setting, the data are divided into
    three parts. The training set fits the models, the validation set
    estimates prediction error for model selection, and the test set assesses
    the generalisation error of the chosen model. The test set should be kept
    \enquote{in a vault} until the end.}
  \defn[§5.3]{goodfellow-2016}{The test set estimates generalisation error
    and must not influence model choices. A validation set, held out from the
    training data, is used to set hyperparameters.}
  \note{In some communities, especially medicine, \enquote{validation} means
    \emph{external} evaluation on new data, which is close to what ML calls
    testing. Check which meaning a document uses.}
  \related{generalization,overfitting,hyperparameter}
\end{entry}

\begin{entry}{deep-learning}{Deep Learning}[DL]
  \defn{lecun-2015}{Representation-learning methods with multiple levels of
    representation, obtained by composing simple non-linear modules that
    each transform the representation at one level into a more abstract
    representation at the next.}
  \defn[ch.~1]{goodfellow-2016}{An approach to AI in which computers learn
    from experience and understand the world as a hierarchy of concepts,
    each defined in relation to simpler ones. The depth of that hierarchy
    gives the field its name.}
  \defn{iso-22989}{An approach to creating rich hierarchical representations
    by training neural networks with many hidden layers.}
  \related{neural-network,machine-learning,transformer}
\end{entry}

\begin{entry}{distribution-shift}{Distribution Shift}[dataset shift, drift]
  \defn{quinonero-2009}{The situation in which the joint distribution of
    inputs and outputs differs between the training and test (or deployment)
    stages. Special cases include covariate shift, where input distributions
    change, and prior probability shift.}
  \defn{gama-2014}{\emph{Concept drift}: a change over time in the
    relationship between input data and the target variable, which makes a
    model learned on past data less accurate.}
  \defn{nist-ai-rmf}{A risk specific to AI: data used to build a system may
    not represent the context of use, and the conditions of deployment may
    change, so that performance and trustworthiness degrade after
    deployment.}
  \note{\enquote{Data drift} usually means a change in the input
    distribution. \enquote{Concept drift} means a change in the input–output
    relationship. Both motivate post-deployment monitoring.}
  \related{generalization,robustness,validity-reliability}
\end{entry}

\begin{entry}{embedding}{Embedding}
  \defn{bengio-2003}{A learned, distributed feature vector that represents a
    discrete symbol, such as a word, in a continuous space, so that similar
    symbols have similar vectors.}
  \defn{mikolov-2013}{A continuous vector representation of words, learned
    from large corpora, in which semantic and syntactic regularities appear
    as vector offsets.}
  \defn[ch.~15]{goodfellow-2016}{Any learned representation that maps
    inputs to points in a vector space where distances reflect a notion of
    similarity useful for a task.}
  \related{deep-learning,retrieval-augmented-generation}
\end{entry}

\begin{entry}{fine-tuning}{Fine-Tuning}
  \defn{howard-ruder-2018}{Adapting a model pre-trained on a large general
    corpus to a target task by continuing to train some or all of its
    parameters on task-specific data.}
  \defn{bommasani-2021}{One form of \emph{adaptation}, in which a foundation
    model's parameters are updated on downstream data. It contrasts with
    prompting, which leaves parameters unchanged.}
  \defn{iso-22989}{Additional training of a pre-trained model on a smaller,
    task- or domain-specific dataset.}
  \related{transfer-learning,foundation-model,rlhf}
\end{entry}

\begin{entry}{generalization}{Generalization}
  \defn[§5.2]{goodfellow-2016}{The ability of a model to perform well on
    previously unobserved inputs, not only on the inputs it was trained on.}
  \defn[ch.~7]{hastie-2009}{Measured by the \emph{generalisation} (test)
    error: the expected prediction error of a model on an independent sample
    from the same distribution as the training data.}
  \defn{iso-22989}{The ability of a machine learning model to make
    predictions that are as accurate on new data as on training data.}
  \note{Classical generalisation assumes that test data come from the same
    distribution as training data. Generalisation under distribution shift is
    a harder, separate problem.}
  \related{overfitting,distribution-shift,datasets}
\end{entry}

\begin{entry}{gradient-descent}{Gradient Descent}
  \defn[§4.3]{goodfellow-2016}{An optimisation method that decreases a
    function by repeatedly moving its parameters a small step in the
    direction of the negative gradient.}
  \defn[§5.2.4]{bishop-2006}{An iterative method for training models in which
    each update moves the weights in the direction of steepest decrease of
    the error function. \emph{Stochastic} or \emph{online} gradient descent
    updates on one data point, or a small batch, at a time.}
  \related{loss-function,neural-network}
\end{entry}

\begin{entry}{hyperparameter}{Hyperparameter}
  \defn[§5.3]{goodfellow-2016}{A setting that controls a learning
    algorithm's behaviour and is not adapted by the algorithm itself, such as
    model capacity or learning rate. Hyperparameters are usually chosen with
    a validation set.}
  \defn{iso-22989}{A characteristic of a machine learning algorithm that
    affects its learning process. It is selected before training rather than
    learned from the training data.}
  \defn{bishop-2006}{A parameter of a prior distribution or of the model
    family, rather than a parameter fitted directly to the data. Examples
    are regularisation coefficients and the number of basis functions.}
  \related{datasets,model}
\end{entry}

\begin{entry}{loss-function}{Loss Function}[cost function, objective]
  \defn[§4.3]{goodfellow-2016}{The function to be minimised during training,
    also called the cost, error, or objective function. It measures how far a
    model's outputs are from the targets.}
  \defn[§2.4]{hastie-2009}{A function $L(Y, f(X))$ that penalises errors in
    prediction. Examples are squared-error loss for regression and 0–1 loss
    for classification.}
  \defn{sutton-barto-2018}{In reinforcement learning, the role of the loss is
    played by a \emph{reward signal} to be maximised rather than an error to
    be minimised.}
  \related{gradient-descent,alignment}
\end{entry}

\begin{entry}{machine-learning}{Machine Learning}[ML]
  \defn*{mitchell-1997}{A computer program is said to learn from experience
    $E$ with respect to some class of tasks $T$ and performance measure $P$,
    if its performance at tasks in $T$, as measured by $P$, improves with
    experience $E$.}
  \defn{samuel-1959}{The field of study that gives computers the ability to
    learn without being explicitly programmed. This wording is widely
    attributed to Samuel but does not appear verbatim in the paper.}
  \defn{iso-22989}{The process of optimising model parameters through
    computational techniques so that the model's behaviour reflects the data
    or experience.}
  \defn[ch.~1]{goodfellow-2016}{The ability of AI systems to acquire their
    own knowledge by extracting patterns from raw data.}
  \note{Mitchell's definition is operational and testable. ISO/IEC 22989
    describes the mechanism, parameter optimisation, and is the one most
    often used in governance documents.}
  \related{artificial-intelligence,supervised-learning,unsupervised-learning,reinforcement-learning}
\end{entry}

\begin{entry}{model}{Model}[ML model]
  \defn{iso-22989}{A mathematical construct that generates an inference or
    prediction from input data. A \emph{machine learning model} is a model
    produced by training on data.}
  \defn[Recital~97]{eu-ai-act}{Distinguished from an \emph{AI system}: a
    model is a component that needs further elements, such as a user
    interface, to become a system.}
  \defn[ch.~1]{bishop-2006}{A parametric function, or a probability
    distribution, whose parameters are fitted to data so that it can predict
    target values for new inputs.}
  \related{ai-system,hyperparameter,model-card}
\end{entry}

\begin{entry}{neural-network}{Neural Network}[artificial neural network]
  \defn{iso-22989}{A network of one or more layers of neurons, connected by
    weighted links with adjustable weights, that takes input data and
    produces an output.}
  \defn[ch.~6]{goodfellow-2016}{A \emph{feedforward network} defines a
    mapping $y = f(x; \theta)$ as a composition of many simpler functions, or
    layers, and learns the parameters $\theta$ that give the best function
    approximation.}
  \defn[ch.~5]{bishop-2006}{A flexible class of parametric non-linear
    functions made of successive linear combinations and non-linear
    activation functions, with parameters fitted by optimisation.}
  \related{deep-learning,transformer,gradient-descent}
\end{entry}

\begin{entry}{overfitting}{Overfitting}
  \defn[§5.2]{goodfellow-2016}{Occurs when the gap between the training
    error and the test error is too large, because the model has memorised
    properties of the training set that do not hold for new data.}
  \defn{iso-22989}{Creating a model that fits the training data too
    precisely and therefore fails to generalise to new data.}
  \defn[§1.1]{bishop-2006}{The phenomenon in which an overly flexible model
    fits noise in the training data, giving small training error but poor
    predictions on new inputs.}
  \related{generalization,datasets}
\end{entry}

\begin{entry}{reinforcement-learning}{Reinforcement Learning}[RL]
  \defn*{sutton-barto-2018}{Reinforcement learning is learning what to
    do---how to map situations to actions---so as to maximize a numerical
    reward signal.}
  \defn{iso-22989}{Learning an optimal sequence of actions to maximise a
    cumulative reward, through trial-and-error interaction with an
    environment.}
  \defn{russell-norvig-2021}{Learning in which an agent learns from a
    series of reinforcements, rewards and punishments, how to act so as to
    maximise expected utility.}
  \related{rlhf,agent,machine-learning}
\end{entry}

\begin{entry}{rlhf}{Reinforcement Learning from Human Feedback}[RLHF]
  \defn{christiano-2017}{Learning a reward function from human judgements
    comparing pairs of agent behaviours, and optimising a policy against that
    learned reward, where writing a reward by hand is impractical.}
  \defn{ouyang-2022}{A procedure for aligning language models with user
    intent. A pre-trained model is fine-tuned on demonstrations, a reward
    model is trained on human rankings of outputs, and the model is then
    optimised against the reward model with reinforcement learning.}
  \defn{ji-2023-alignment}{A family of \emph{learning from feedback}
    techniques used for forward alignment, in which human preference signals
    steer model behaviour toward intended values.}
  \related{alignment,fine-tuning,reinforcement-learning}
\end{entry}

\begin{entry}{self-supervised-learning}{Self-Supervised Learning}[SSL]
  \defn{liu-2021-ssl}{Learning in which supervisory signals are derived
    automatically from the data itself, for example by predicting masked or
    future parts of the input. It needs no human-annotated labels.}
  \defn{bommasani-2021}{The training method behind foundation models. A
    model is trained on unlabelled data with an objective derived from the
    data, such as predicting the next word.}
  \defn{iso-22989}{Machine learning that uses unlabelled data and
    automatically generated labels. It is sometimes treated as a subset of
    unsupervised learning.}
  \related{foundation-model,unsupervised-learning}
\end{entry}

\begin{entry}{supervised-learning}{Supervised Learning}
  \defn{iso-22989}{Machine learning that uses only labelled data during
    training.}
  \defn[ch.~2]{hastie-2009}{Learning to predict an outcome from inputs using a
    training set of observed input–outcome pairs, where the outcome acts as
    a \enquote{teacher} guiding learning.}
  \defn[§5.1.3]{goodfellow-2016}{Learning from a dataset in which each
    example is associated with a label or target, typically in order to
    estimate $p(y \mid x)$.}
  \related{unsupervised-learning,self-supervised-learning,datasets}
\end{entry}

\begin{entry}{transfer-learning}{Transfer Learning}
  \defn{pan-yang-2010}{Improving learning of a target task in a target
    domain by using knowledge from a different source domain or task,
    particularly when target data are scarce.}
  \defn[§15.2]{goodfellow-2016}{Using what has been learned in one setting,
    such as one distribution or task, to improve generalisation in another
    setting.}
  \defn{bommasani-2021}{The mechanism that makes foundation models possible:
    knowledge learned during pre-training on a surrogate task is transferred
    to downstream tasks through adaptation.}
  \related{fine-tuning,foundation-model}
\end{entry}

\begin{entry}{transformer}{Transformer}
  \defn*{vaswani-2017}{A new simple network architecture, the Transformer,
    based solely on attention mechanisms, dispensing with recurrence and
    convolutions entirely.}
  \defn{zhao-2023}{The standard backbone of large language models. It uses
    stacked multi-head self-attention and feed-forward layers, and its
    parallelism allows training on very large corpora.}
  \defn{bommasani-2021}{An architecture whose scalability and ability to
    handle multiple modalities made it the dominant basis of foundation
    models.}
  \related{large-language-model,neural-network,deep-learning}
\end{entry}

\begin{entry}{unsupervised-learning}{Unsupervised Learning}
  \defn{iso-22989}{Machine learning that uses only unlabelled data during
    training.}
  \defn[ch.~14]{hastie-2009}{Learning without a \enquote{teacher}: inferring
    properties of the probability density of the inputs directly, for
    example through clustering, density estimation, or dimensionality
    reduction.}
  \defn[§5.1.3]{goodfellow-2016}{Learning useful properties of the structure
    of a dataset, often the probability distribution that generated it,
    from examples without labels.}
  \related{supervised-learning,self-supervised-learning}
\end{entry}

% Trustworthy AI terms, in alphabetical order.
\domain{Trustworthy AI}

\begin{entry}{accountability}{Accountability}
  \defn{hleg-2019}{A requirement for trustworthy AI: mechanisms are in place
    to ensure responsibility and accountability for AI systems and their
    outcomes, before and after development, deployment, and use. These
    include auditability, minimising and reporting negative impacts,
    trade-offs, and redress.}
  \defn[§3.4]{nist-ai-rmf}{Accountability presupposes transparency. It
    concerns who answers for an AI system's outcomes, and it depends on
    organisational practices and governing structures for reducing harms,
    such as risk management and clear roles and responsibilities.}
  \defn{bovens-2007}{In public governance, a relationship between an actor
    and a forum in which the actor must explain and justify its conduct, the
    forum can pose questions and pass judgement, and the actor may face
    consequences.}
  \note{Bovens's relational definition clarifies that accountability needs
    an identifiable party who must answer, and someone they must answer to.
    Transparency alone does not provide either.}
  \related{transparency,auditability,ai-governance}
\end{entry}

\begin{entry}{adversarial-example}{Adversarial Example}
  \defn{szegedy-2014}{An input formed by applying a small, often
    imperceptible, perturbation to a correctly classified example, chosen so
    that a neural network misclassifies it.}
  \defn{goodfellow-2015}{Inputs formed by intentionally adding worst-case
    perturbations. The authors argue that vulnerability to them comes mainly
    from the models' largely linear behaviour in high-dimensional spaces,
    not from non-linearity or overfitting.}
  \defn{nist-ai-100-2}{The product of an \emph{evasion attack}: at
    deployment time, an adversary modifies a test sample to change the
    model's prediction to one the attacker chooses.}
  \related{robustness,data-poisoning,security-resilience}
\end{entry}

\begin{entry}{ai-governance}{AI Governance}
  \defn{mantymaki-2022}{A system of rules, practices, processes, and
    technological tools used to ensure that an organisation's use of AI
    technologies is aligned with its strategies, objectives, and values,
    fulfils legal requirements, and meets principles of ethical AI.}
  \defn{iso-42001}{Realised through an \emph{AI management system}: the
    interrelated policies, objectives, and processes an organisation
    establishes for the responsible development, provision, or use of AI
    systems.}
  \defn[§5.1]{nist-ai-rmf}{The \textsc{Govern} function: a culture of risk
    management is cultivated and present across the organisation. It cuts
    across and informs the \textsc{Map}, \textsc{Measure}, and
    \textsc{Manage} functions.}
  \related{accountability,ai-risk,auditability}
\end{entry}

\begin{entry}{ai-risk}{AI Risk}
  \defn*[Art.~3(2)]{eu-ai-act}{The combination of the probability of an
    occurrence of harm and the severity of that harm.}
  \defn[§1.1]{nist-ai-rmf}{A composite measure of an event's probability of
    occurring and the magnitude or degree of the consequences. AI risks can
    affect individuals, groups, organisations, communities, society, the
    environment, and the planet.}
  \defn{iso-31000}{The effect of uncertainty on objectives, where an effect
    is a deviation from the expected and can be positive, negative, or
    both.}
  \note{ISO 31000 counts positive deviations (opportunities) as risk. The EU
    AI Act and NIST AI RMF concentrate on harm. The EU AI Act classifies
    systems into risk tiers (prohibited, high-risk, limited, and minimal),
    and each tier carries different obligations.}
  \related{ai-governance,safety,red-teaming}
\end{entry}

\begin{entry}{alignment}{Alignment}[AI alignment]
  \defn{ji-2023-alignment}{The goal of making AI systems behave in line with
    human intentions and values. It is characterised by the RICE objectives:
    robustness, interpretability, controllability, and ethicality.}
  \defn{gabriel-2020}{The normative and technical problem of making AI act
    in accordance with appropriate values. The problem includes choosing
    whether a system should follow instructions, intentions, revealed
    preferences, ideal preferences, interests, or values.}
  \defn{russell-2019}{Ensuring that machines pursue objectives that are
    beneficial to humans. The proposed route is for machines to be uncertain
    about human preferences and to learn them from behaviour, rather than to
    optimise a fixed, possibly misspecified objective.}
  \related{rlhf,safety,human-oversight}
\end{entry}

\begin{entry}{auditability}{Auditability}[algorithmic audit]
  \defn{hleg-2019}{The ability to assess an AI system's algorithms, data, and
    design processes, by internal or external auditors, with the resulting
    reports available. It is especially important for systems affecting
    fundamental rights.}
  \defn{raji-2020}{Supported by an internal audit process, SMACTR (scoping,
    mapping, artifact collection, testing, and reflection), that produces
    documentation throughout development. The documentation lets an
    organisation check a system against its declared ethical expectations
    before deployment.}
  \defn{iso-42001}{Achieved through internal audits of the AI management
    system at planned intervals, which check conformity with the
    organisation's requirements and the standard.}
  \related{accountability,transparency,model-card}
\end{entry}

\begin{entry}{bias}{Bias}[AI bias]
  \defn{nist-sp-1270}{Bias in AI falls into three categories.
    \emph{Systemic} biases come from institutional practices.
    \emph{Statistical and computational} biases come from unrepresentative
    data or modelling choices. \emph{Human-cognitive} biases shape how
    people design, interpret, and use AI systems.}
  \defn{friedman-nissenbaum-1996}{Computer systems that \emph{systematically
    and unfairly discriminate} against certain individuals or groups in
    favour of others. Bias may be pre-existing, technical, or emergent.}
  \defn{mehrabi-2021}{Systematic errors or skews at any stage of the ML
    pipeline, such as data collection, sampling, measurement, algorithm
    design, or user interaction, that can lead to unfair outcomes.}
  \defn{hastie-2009}{In statistics, the difference between an estimator's
    expected value and the true value. This is a technical property with no
    moral meaning, and is often traded off against variance.}
  \note{The statistical meaning is neutral and often intentional, as in
    regularisation. The socio-technical meanings are about unjust outcomes.
    Say which meaning is intended.}
  \related{fairness,datasets}
\end{entry}

\begin{entry}{data-poisoning}{Data Poisoning}
  \defn{nist-ai-100-2}{An attack at training time in which an adversary
    controls part of the training data, or the labels, to degrade model
    performance in general (availability poisoning) or to cause targeted or
    backdoor behaviour (integrity poisoning).}
  \defn{biggio-2012}{Injecting specially crafted training points to raise a
    learning algorithm's test error. The paper demonstrates this against
    support vector machines by computing the points with gradient ascent.}
  \defn{owasp-llm-2025}{For large language models: manipulation of
    pre-training, fine-tuning, or embedding data to introduce
    vulnerabilities, backdoors, or biases that compromise security,
    performance, or ethical behaviour.}
  \related{adversarial-example,security-resilience,datasets}
\end{entry}

\begin{entry}{differential-privacy}{Differential Privacy}[DP]
  \defn{dwork-roth-2014}{A mathematical guarantee that the output of an
    analysis is almost equally likely whether or not any single individual's
    data are included. A mechanism $M$ is $(\varepsilon,\delta)$-differentially
    private if, for all neighbouring datasets $D, D'$ and all sets of outputs
    $S$, $\Pr[M(D)\in S] \le e^{\varepsilon}\Pr[M(D')\in S] + \delta$.}
  \defn[§3.6]{nist-ai-rmf}{One of several privacy-enhancing technologies
    that can support privacy-enhanced AI. It may trade off against accuracy,
    and therefore against fairness for small groups.}
  \related{privacy}
\end{entry}

\begin{entry}{explainability}{Explainability}[XAI]
  \defn[§3.5]{nist-ai-rmf}{A representation of the mechanisms underlying an
    AI system's operation. It answers the question of \emph{how} a decision
    was made.}
  \defn{nist-ir-8312}{Characterised by four principles. A system gives
    evidence or reasons for its outputs (\emph{explanation}). The
    explanations are understandable to their audience (\emph{meaningful}).
    They correctly reflect the process that produced the output
    (\emph{explanation accuracy}). The system operates only within its
    designed conditions and confidence (\emph{knowledge limits}).}
  \defn{barredo-2020}{Given an audience, an explainable AI is one that
    produces details or reasons to make its functioning clear or easy to
    understand.}
  \defn{hleg-2019}{Part of transparency: the ability to explain both the
    technical processes of an AI system and the related human decisions,
    adapted to the expertise of the stakeholder concerned.}
  \note{Every definition makes explainability depend on its \emph{audience}.
    NIST separates it from interpretability, but many papers use the two
    words interchangeably.}
  \related{interpretability,transparency}
\end{entry}

\begin{entry}{fairness}{Fairness}[algorithmic fairness]
  \defn[§3.7]{nist-ai-rmf}{The characteristic \enquote{fair --- with harmful
    bias managed}: concerns for equality and equity, addressed by managing
    harmful bias and discrimination. Standards of fairness are complex and
    differ across cultures and applications.}
  \defn{dwork-2012}{\emph{Individual fairness}: similar individuals should be
    treated similarly, formalised as a Lipschitz condition on the classifier
    under a task-specific similarity metric.}
  \defn{hardt-2016}{\emph{Equalised odds}: a predictor is fair if its
    predictions are independent of a protected attribute conditional on the
    true outcome. That is, true- and false-positive rates are equal across
    groups. \emph{Equal opportunity} relaxes this to equal true-positive
    rates.}
  \defn{hleg-2019}{A requirement named \emph{diversity, non-discrimination
    and fairness}. It covers avoiding unfair bias, accessibility and
    universal design, and involving stakeholders throughout the lifecycle.}
  \note{Group criteria, such as demographic parity and equalised odds, and
    individual criteria cannot in general all be satisfied at once. Choosing
    among them is a normative decision that should be documented.}
  \related{bias,accountability}
\end{entry}

\begin{entry}{human-oversight}{Human Oversight}
  \defn[Art.~14]{eu-ai-act}{High-risk AI systems must be designed and
    developed so that natural persons can effectively oversee them while
    they are in use. The aim is to prevent or minimise risks to health,
    safety, or fundamental rights, including by being able to understand,
    monitor, override, or stop the system.}
  \defn{hleg-2019}{Part of the requirement \emph{human agency and
    oversight}. It may be achieved through human-in-the-loop,
    human-on-the-loop, or human-in-command approaches, which differ in how
    closely a human is involved.}
  \defn{nist-ai-rmf}{Human roles and responsibilities in decision-making and
    in overseeing AI systems should be clearly defined and differentiated,
    taking account of human cognitive biases such as automation bias.}
  \note{Oversight is only effective if the overseer has the competence,
    authority, time, and information to intervene. Without these, a human
    reviewer may end up bearing responsibility for outcomes they could not
    actually control.}
  \related{human-in-the-loop,autonomy,accountability}
\end{entry}

\begin{entry}{human-in-the-loop}{Human-in-the-Loop}[HITL]
  \defn{hleg-2019}{The capability for human intervention in every decision
    cycle of the system. It contrasts with \emph{human-on-the-loop}, where
    humans intervene during design and monitor operation, and
    \emph{human-in-command}, where humans oversee overall activity and
    decide when and how to use the system.}
  \defn{mosqueira-rey-2023}{In machine learning, approaches in which humans
    take part in the learning process, for example through active learning,
    interactive ML, or machine teaching, so as to improve model accuracy or
    efficiency.}
  \note{The governance meaning concerns human control over decisions. The
    ML meaning concerns humans contributing to training.}
  \related{human-oversight,rlhf}
\end{entry}

\begin{entry}{interpretability}{Interpretability}
  \defn*{doshi-velez-2017}{The ability to explain or to present in
    understandable terms to a human.}
  \defn{miller-2019}{The degree to which a human can understand the cause of
    a decision. Miller uses it interchangeably with \emph{explainability}.}
  \defn[§3.5]{nist-ai-rmf}{The meaning of an AI system's output in the
    context of its designed functional purpose. It answers the question of
    \emph{why} a decision was made and what it means to the user.}
  \related{explainability,transparency}
\end{entry}

\begin{entry}{model-card}{Model Card}[datasheet]
  \defn{mitchell-2019}{A short document accompanying a trained model that
    gives benchmarked evaluation across relevant conditions, such as
    demographic groups, as well as intended use, out-of-scope uses,
    evaluation procedures, and ethical considerations.}
  \defn{gebru-2021}{The companion practice for data, called a
    \emph{datasheet for datasets}. It documents a dataset's motivation,
    composition, collection process, preprocessing, uses, distribution, and
    maintenance.}
  \defn[Art.~11]{eu-ai-act}{Related regulatory concept: the \emph{technical
    documentation} that providers of high-risk AI systems must draw up
    before placing them on the market and keep up to date.}
  \related{transparency,auditability,datasets}
\end{entry}

\begin{entry}{privacy}{Privacy}[privacy-enhanced AI]
  \defn[§3.6]{nist-ai-rmf}{The norms and practices that help to safeguard
    human autonomy, identity, and dignity. These include freedom from
    intrusion, limits on observation, and individuals' agency to consent to
    disclosure or control of facets of their identities.}
  \defn{hleg-2019}{The requirement \emph{privacy and data governance}: respect
    for privacy, the quality and integrity of data, and legitimate access to
    data, throughout the system's lifecycle.}
  \defn{nist-ai-100-2}{The target of \emph{privacy attacks}, such as
    membership inference, data reconstruction, and model extraction, which
    aim to recover information about training data or the model.}
  \related{differential-privacy,security-resilience}
\end{entry}

\begin{entry}{prompt-injection}{Prompt Injection}
  \defn{nist-ai-100-2}{An attack in which an adversary inserts instructions
    into a generative AI system's input. In \emph{direct} prompt injection
    the attacker is the user. In \emph{indirect} prompt injection the
    instructions arrive through external content, such as a web page or
    document, that the system consumes.}
  \defn{greshake-2023}{Indirect prompt injection: planting prompts in data
    likely to be retrieved by an LLM-integrated application, which blurs the
    line between data and instructions and allows the attacker to control
    the application remotely.}
  \defn{owasp-llm-2025}{A vulnerability, ranked first among risks to LLM
    applications, in which user prompts or external content alter the
    model's behaviour or output in unintended ways. Outcomes include data
    disclosure and unauthorised actions.}
  \note{Prompt injection is distinct from \emph{jailbreaking}, which targets
    the model's safety training. Prompt injection targets the application
    built around the model.}
  \related{prompt,agentic-ai,retrieval-augmented-generation,security-resilience}
\end{entry}

\begin{entry}{red-teaming}{Red Teaming}
  \defn{ganguli-2022}{Adversarially probing a language model, by human or
    automated means, to discover and measure harmful outputs, so that the
    harms can be analysed and reduced.}
  \defn{nist-ai-600-1}{A structured testing effort, often by teams outside
    the development group, to find flaws and vulnerabilities in a generative
    AI system, such as harmful or discriminatory outputs, unforeseen
    behaviour, or potential misuse.}
  \defn[Art.~55]{eu-ai-act}{Adversarial testing that providers of
    general-purpose AI models with systemic risk must carry out and
    document, in order to identify and mitigate systemic risks.}
  \related{ai-risk,safety,security-resilience}
\end{entry}

\begin{entry}{robustness}{Robustness}
  \defn{iso-22989}{The ability of a system to maintain its level of
    performance under any circumstances, including external interference
    and harsh environmental conditions.}
  \defn{hleg-2019}{Named \emph{technical robustness and safety}: AI systems
    are resilient to attack, have fall-back plans, and are accurate,
    reliable, and reproducible, so that unintentional harm is minimised.}
  \defn[Art.~15]{eu-ai-act}{High-risk AI systems shall be as resilient as
    possible to errors, faults, or inconsistencies within the system or its
    environment, for example through technical redundancy and fail-safe
    plans.}
  \defn{madry-2018}{\emph{Adversarial robustness}: low worst-case loss under
    any perturbation within an allowed set, formalised as a min–max
    (saddle-point) optimisation problem.}
  \related{adversarial-example,validity-reliability,distribution-shift}
\end{entry}

\begin{entry}{safety}{Safety}[AI safety]
  \defn[§3.2]{nist-ai-rmf}{AI systems should not, under defined conditions,
    lead to a state in which human life, health, property, or the
    environment is endangered.}
  \defn{amodei-2016}{The problem of \emph{accidents} in ML systems:
    unintended and harmful behaviour arising from poor design. Examples are
    negative side effects, reward hacking, unsafe exploration, and
    brittleness under distributional shift.}
  \defn{hleg-2019}{Part of technical robustness: AI systems are developed
    with a preventive approach to risk, behave reliably as intended, and
    minimise unintentional and unexpected harm.}
  \note{\enquote{AI safety} also names a research field concerned with
    risks from advanced AI, including misalignment. In product and
    engineering contexts it usually means freedom from unacceptable risk of
    physical or psychological harm.}
  \related{alignment,ai-risk,robustness}
\end{entry}

\begin{entry}{security-resilience}{Security and Resilience}
  \defn[§3.3]{nist-ai-rmf}{An AI system is \emph{resilient} if it can
    withstand unexpected adverse events or changes in its environment or use,
    or degrade safely and gracefully. It is \emph{secure} if it maintains
    confidentiality, integrity, and availability through mechanisms that
    prevent unauthorised access and use.}
  \defn[Art.~15]{eu-ai-act}{Cybersecurity: high-risk AI systems shall be
    resilient against attempts by unauthorised third parties to alter their
    use, outputs, or performance by exploiting vulnerabilities. Examples
    include data poisoning, model poisoning, adversarial examples, and
    confidentiality attacks.}
  \defn{nist-ai-100-2}{Adversarial machine learning classifies attacks by
    the attacker's goals (availability, integrity, privacy, and misuse),
    capabilities, knowledge, and the stage of the ML lifecycle they target.}
  \related{adversarial-example,data-poisoning,prompt-injection,robustness}
\end{entry}

\begin{entry}{transparency}{Transparency}
  \defn[§3.4]{nist-ai-rmf}{The extent to which information about an AI
    system and its outputs is available to the individuals who interact with
    it, whether or not they know they are doing so.}
  \defn{hleg-2019}{A requirement covering \emph{traceability} of data and
    processes, \emph{explainability}, and \emph{communication}. People are
    told when they are interacting with an AI system, and they are told its
    capabilities and limitations.}
  \defn[Art.~13]{eu-ai-act}{High-risk AI systems shall be designed so that
    their operation is sufficiently transparent to enable deployers to
    interpret the system's output and use it appropriately. The systems are
    accompanied by instructions for use.}
  \defn[Art.~50]{eu-ai-act}{Transparency obligations toward the people
    affected: for example, disclosing that they are interacting with an AI
    system, and marking synthetic or deep-fake content.}
  \related{explainability,accountability,model-card}
\end{entry}

\begin{entry}{trustworthy-ai}{Trustworthy AI}
  \defn{hleg-2019}{AI that is \emph{lawful}, \emph{ethical}, and
    \emph{robust}. It is realised through seven requirements: human agency
    and oversight; technical robustness and safety; privacy and data
    governance; transparency; diversity, non-discrimination and fairness;
    societal and environmental well-being; and accountability.}
  \defn[§3]{nist-ai-rmf}{AI that has these characteristics: valid and
    reliable; safe; secure and resilient; accountable and transparent;
    explainable and interpretable; privacy-enhanced; and fair with harmful
    bias managed. Validity and reliability are the base for the others, and
    the characteristics must be balanced against each other in context.}
  \defn{iso-24028}{Rests on \emph{trustworthiness}: the ability to meet
    stakeholders' expectations in a verifiable way. For AI this covers
    characteristics such as reliability, availability, resilience, security,
    privacy, safety, accountability, transparency, integrity, authenticity,
    quality, and usability.}
  \defn{oecd-2024}{AI that respects the OECD values-based principles:
    inclusive growth, sustainable development and well-being; human rights
    and democratic values, including fairness and privacy; transparency and
    explainability; robustness, security and safety; and accountability.}
  \note{The frameworks overlap heavily but are organised differently. The
    EU HLEG starts from ethics and fundamental rights, NIST from risk
    management, and ISO/IEC from verifiable stakeholder expectations.}
  \related{accountability,explainability,fairness,privacy,robustness,safety,transparency,validity-reliability}
\end{entry}

\begin{entry}{validity-reliability}{Validity and Reliability}
  \defn[§3.1]{nist-ai-rmf}{\emph{Validation} is confirmation, through
    objective evidence, that the requirements for a specific intended use
    have been fulfilled. \emph{Reliability} is the ability of an item to
    perform as required, without failure, for a given time interval under
    given conditions. Accuracy and robustness contribute to both.}
  \defn{iso-22989}{Reliability is the property of consistent intended
    behaviour and results.}
  \defn{hleg-2019}{Part of technical robustness: a reliable system works
    properly across a range of inputs and situations, and reproducibility
    means it behaves the same way when an experiment is repeated under the
    same conditions.}
  \note{NIST treats validity and reliability as the foundation of
    trustworthiness. A system that is not valid for its context cannot be
    made trustworthy by adding the other characteristics.}
  \related{robustness,trustworthy-ai,distribution-shift}
\end{entry}

\input{entries/notes}
\input{entries/acronyms}

\clearpage
\markboth{Sources}{Sources}
\printbibliography[heading=bibintoc,title=Sources]

\clearpage
\markboth{Index}{Index}
\printindex

\end{document}
